\documentclass[11pt,a4paper]{article}
\usepackage[T1]{fontenc}
\usepackage[margin=26mm]{geometry}
\usepackage{amsmath,amssymb,amsthm,mathtools}
\usepackage{enumitem}
\usepackage[colorlinks=true,linkcolor=blue,citecolor=blue,urlcolor=blue]{hyperref}
\usepackage{microtype}
\setlength{\parskip}{0.35em}
\setlist{nosep}
\numberwithin{equation}{section}
\newtheorem{theorem}{Theorem}[section]
\newtheorem{lemma}[theorem]{Lemma}
\newtheorem{proposition}[theorem]{Proposition}
\newtheorem{corollary}[theorem]{Corollary}
\theoremstyle{definition}
\newtheorem{definition}[theorem]{Definition}
\theoremstyle{remark}
\newtheorem{remark}[theorem]{Remark}
\newcommand{\R}{\mathbb R}
\newcommand{\dd}{\,\mathrm d}
\newcommand{\dx}{\dd x}
\newcommand{\dy}{\dd y}
\newcommand{\dz}{\dd z}
\newcommand{\avint}{\mathop{\fint}}
\newcommand{\norm}[1]{\left\lVert#1\right\rVert}
\newcommand{\E}{\mathcal E}
\newcommand{\Th}{\Theta}
\newcommand{\osc}{\operatorname{osc}}
\newcommand{\supp}{\operatorname{supp}}
\newcommand{\dist}{\operatorname{dist}}
\newcommand{\PV}{\operatorname{p.v.}}
\title{From H\"older representatives to quantitative smooth regularity\\
\large A complete interior upgrade for energy-minimizing harmonic maps}
\author{}
\date{}
\begin{document}
\maketitle

\begin{abstract}
We prove that the quantitative small-energy H\"older representative theorem
for locally energy-minimizing harmonic maps implies smoothness and estimates
for every positive-order ambient derivative. The energy normalization is
\(\int |Du|_{\mathrm{HS}}^2\), without a factor of one half. A single energy
threshold works for all derivative orders. The proof first treats the weak
quadratic equation by harmonic replacement, improves energy growth, derives
gradient Campanato decay, and identifies the continuous weak gradient with
the classical derivative. A constant-coefficient Poisson Schauder estimate is
proved using the fundamental solution and its cancellation, rather than
assumed. Differentiating the smooth coefficient equation then gives all
higher estimates on nested balls. The endpoint is interior regularity on
Euclidean domains; no claim about boundary regularity is made.
\end{abstract}

\section{Local data, representatives, and the quantitative target}
\label{sec:local}

Fix \(n\ge2\) and a nonempty compact smooth embedded submanifold
\(N\subset\R^q\), without boundary. The metric on \(N\) is induced by the
ambient Euclidean inner product. Balls are open unless explicitly stated
otherwise. The norm of \(Du\) in energy integrals is the Hilbert--Schmidt
norm; the norms of \(D^j u\), \(j\ge1\), are multilinear operator norms.
These norms are equivalent to their finite coordinate norms, with constants
depending only on \(n,q,j\). For a continuous function \(v\) on \(U\), write
\[
 [v]_{\eta;U}=\sup_{x\ne y\in U}
 \frac{|v(x)-v(y)|}{|x-y|^\eta},\qquad 0<\eta<1.
\]
For functions with \(m\) continuous derivatives define
\[
 \norm{v}_{C^{m,\eta}(U)}
 =\sum_{j=0}^m\sup_U\norm{D^jv}+[D^mv]_{\eta;U}.
\]
All uses of these norms below concern relatively compact balls, where the
displayed quantities are finite after the corresponding estimate is proved.

An ambient Sobolev map is \(u\in W^{1,2}_{\mathrm{loc}}(\Omega,\R^q)\)
with its displayed distributional gradient \(Du\), and \(u\in N\) almost
everywhere. Local minimality means
\[
 \int_{B_r(b)}|Du|^2\le\int_{B_r(b)}|Dv|^2
 \quad\text{if }\overline{B_r(b)}\subset\Omega,
 \quad v\in W^{1,2}(B_r(b),N),\quad v-u\in W^{1,2}_0(B_r(b),\R^q).
\]
Here \(W^{1,2}_0\) is the closure of smooth compactly supported maps in
the function-and-gradient \(L^2\) norm. Define
\[
 \E(u;B)=\int_B |Du|^2,\qquad
 \Th_u(a,R)=R^{2-n}\E(u;B_R(a)).
\]

\begin{theorem}[Quantitative smooth small-energy regularity]
\label{thm:main}
There are \(\varepsilon_*>0\), \(\gamma\in(0,1)\), and numbers
\(C_j<\infty\), \(j\ge1\), depending only on \(n,N\) (and on \(j\)
for \(C_j\)), with the following property. Suppose \(u\) is locally energy
minimizing, \(\overline{B_R(a)}\subset\Omega\), and
\(\Th_u(a,R)\le\varepsilon_*\). The unique continuous representative
\(w\) on \(B_{R/2}(a)\) takes values in \(N\), is smooth, and satisfies
\begin{equation}\label{eq:main}
 \begin{aligned}
 R^j\sup_{B_{R/2}(a)}\norm{D^jw}
 &\le C_j\Th_u(a,R)^{1/2},\\
 R^{j+\gamma}[D^jw]_{\gamma;B_{R/2}(a)}
 &\le C_j\Th_u(a,R)^{1/2},\qquad j\ge1.
 \end{aligned}
\end{equation}
It also satisfies
\begin{equation}\label{eq:zero}
 \sup_{B_{R/2}(a)}|w-w(a)|\le C_0\Th_u(a,R)^{1/2}
\end{equation}
for a constant depending only on \(n,N\). The sum of the two scaled
positive-order quantities is consequently bounded by
\(2C_j\Th_u(a,R)^{1/2}\). The threshold is selected before
the map, domain, center, radius, and derivative order. Uniqueness concerns
the representative on the indicated ball.
\end{theorem}

\begin{remark}
The zero-order bound concerns \(w-w(a)\). A nonzero constant map has zero
energy, so a bound for the uncentered full \(C^j\) norm by
\(C_j\Th^{1/2}\) would be false. Ambient total functions may be used to
represent maps in a proof assistant, but all equality, smoothness, and
derivative assertions in this theorem are restricted to the open ball.
\end{remark}

We use the already established H\"older theorem \cite{Holder}: there are
\(\varepsilon_H>0\), \(\alpha\in(0,1)\), and \(C_H\), depending only on
\(n,N\), such that small normalized energy on a compactly contained ball
gives a unique \(N\)-valued representative on its half-ball, with
\[
 [w]_{\alpha;B_{s/2}(b)}
 \le C_Hs^{-\alpha}\Th_u(b,s)^{1/2}.
\]
The other inputs from that development are the genuine weak equation,
smooth nearest-point projection and its smooth Euclidean extension,
Euclidean Sobolev mollification, the smooth-data Poisson solution on a ball,
and its signed Green and energy-orthogonality identities. They are
foundational inputs, not assumptions about the higher regularity of \(u\).
The new elliptic regularity arguments are proved in this note.

\begin{lemma}[Uniform H\"older data on an enlarged inner ball]
\label{lem:localholder}
Suppose the unit-ball energy is \(E\), and
\(E\le4^{2-n}\varepsilon_H\). There is an \(N\)-valued continuous
representative on \(B_{3/4}\) such that
\begin{equation}\label{eq:holderinput}
 [w]_{\alpha;B_{3/4}}\le L\sqrt E,
 \qquad \int_{B_{3/4}}|Dw|^2\le E,
\end{equation}
where \(Dw\) in the second expression initially means the same weak
gradient as \(Du\), and \(L\) depends only on \(n,N\).
\end{lemma}
\begin{proof}
For \(x\in B_{3/4}\), the closed ball of radius \(1/4\) centered at \(x\)
lies in \(B_1\), and its normalized energy is at most \(4^{n-2}E\).
Apply the H\"older theorem on this ball. Its representative on \(B_{1/8}(x)\)
has seminorm at most
\(L_0\sqrt E\), where \(L_0=C_H4^{\alpha+(n-2)/2}\).
Representatives agree on overlaps: their difference is continuous and zero
almost everywhere on an open set, and every nonempty Euclidean open ball
has positive measure. Thus they glue. A countable subcover exists by second
countability, so the glued function agrees almost everywhere with \(u\);
no uncountable union of exceptional null sets is used.
For \(x,y\in B_{3/4}\), divide their segment into 24 equal parts. Every
part has length less than \(1/16\), and lies in a ball \(B_{1/8}\) centered
at its initial point. Summing the local estimates gives
\(L=24^{1-\alpha}L_0\). The equality of weak gradients follows from
distributional almost-everywhere invariance and uniqueness; restriction
of the energy gives the second bound.
\end{proof}

Let \(\Pi\) be the nearest-point projection on a fixed tubular neighborhood.
Choose a smooth Euclidean extension \(Q:\R^q\to\R^q\) agreeing with
\(\Pi\) near \(N\), and put
\begin{equation}\label{eq:coefficient}
 A(z)[v,t]=D^2Q(z)[v,t].
\end{equation}
At target points this equals the canonical projection Hessian, since
the two maps agree on a neighborhood. The established outer-variation
equation gives, on \(B_{3/4}\),
\begin{equation}\label{eq:weak}
 \Delta w=F,\qquad
 F=\sum_{i=1}^n A(w)[D_iw,D_iw],\qquad |F|\le K|Dw|^2
\end{equation}
almost everywhere, with \(K<\infty\) depending only on \(N\). Precisely,
for every smooth compactly supported vector test \(\psi\),
\[
 \int Dw:D\psi=-\int F\cdot\psi.
\]
The right side is integrable because \(F\in L^1_{\mathrm{loc}}\) and the
test is bounded. This sign convention is used in every Green identity.

\section{Harmonic replacement and the weak quadratic equation}
\label{sec:comparison}

We first record the harmonic estimates needed below, including their scaling.
For a smooth harmonic vector map \(h\) on \(B_r(x)\),
\begin{align}
 \int_{B_{\theta r}(x)}|Dh|^2
 &\le c_n\theta^n\int_{B_r(x)}|Dh|^2,
 \label{eq:henergy}\\
 \int_{B_{\theta r}(x)}|Dh-(Dh)_{x,\theta r}|^2
 &\le c_n\theta^{n+2}\int_{B_r(x)}|Dh-(Dh)_{x,r}|^2,
 \label{eq:hexcess}
\end{align}
for \(0<\theta\le1/2\). Here a subscript \((x,r)\) denotes a ball mean.
To see all derivative estimates used in these inequalities, choose a smooth
radial kernel \(\zeta\) supported in \(B_1\), of integral one. The ball
mean-value formula, integrated against its radial profile, gives
\(h=\zeta_s*h\) whenever the sampling ball is contained in the harmonic
domain. Differentiating the kernel and using Cauchy--Schwarz gives
\begin{equation}\label{eq:hderivatives}
 \sup_{B_{r/2}(x)}\norm{D^jh}
 \le c_{n,j}r^{-j-n/2}\norm{h}_{L^2(B_r(x))}.
\end{equation}
The same argument with an \(L^1\) or supremum bound gives the corresponding
interior estimates. Apply the mean-value formula to derivative columns for
\eqref{eq:henergy}. For \eqref{eq:hexcess}, apply
\eqref{eq:hderivatives} with \(j=1\) to the harmonic field
\(Dh-(Dh)_{x,r}\), then integrate its Lipschitz oscillation on the smaller
ball. All finite-coordinate/operator-norm conversions enter \(c_n\).

\begin{lemma}[Harmonic replacement for a continuous Sobolev datum]
\label{lem:replacement}
Let \(\overline{B_r(x)}\subset B_{3/4}\). There is
\(h\in W^{1,2}(B_r(x),\R^q)\cap C(\overline{B_r(x)},\R^q)\),
smooth and harmonic in the open ball, with boundary value \(w\), such that
\(w-h\in W^{1,2}_0(B_r(x),\R^q)\),
\begin{equation}\label{eq:orthogonal}
 \int_{B_r(x)}|Dw|^2
 =\int_{B_r(x)}|Dh|^2+\int_{B_r(x)}|Dw-Dh|^2,
\end{equation}
and, if \([w]_{\alpha;B_{3/4}}\le H\),
\begin{equation}\label{eq:error}
 \int_{B_r(x)}|Dw-Dh|^2
 \le 2KHr^\alpha\int_{B_r(x)}|Dw|^2.
\end{equation}
\end{lemma}
\begin{proof}
Fix a slightly larger compactly contained ball. Mollify the continuous
Sobolev map in this larger ball by a nonnegative smooth unit-mass kernel
supported at radius \(t\), obtaining \(w_t\). Local convolution gives
uniform convergence on the closed ball, strong function-and-gradient
\(L^2\) convergence, and
\begin{equation}\label{eq:mollpoisson}
 \Delta w_t=F_t,\qquad
 \int_{B_r(x)}|F_t|\le K\int_{B_{r+t}(x)}|Dw|^2.
\end{equation}
For the Laplacian identity, gradient commutation first gives
\(D_iw_t(x)=\int\varphi_t(x-y)D_iw(y)\dy\).
Differentiate the smooth kernel in \(x_i\) and sum over \(i\).
In the weak Poisson equation use the vector test
\(\psi(y)=\varphi_t(x-y)e_j\). Its derivative in \(y_i\) is
\(-D_i\varphi_t(x-y)e_j\), so the two negative signs in the signed test
identity cancel and give the \(j\)-th component of
\(\Delta w_t=\varphi_t*F\). The kernel support is compactly contained
in the domain, and local integrability supplies all these pairings.
Jensen's inequality and Fubini give the displayed forcing integral bound.
The other assertions follow by the weak-gradient test identities and
translation continuity in \(L^2\). A cutoff in the larger ball can make the convolutions
global without changing them on the closed ball under consideration.

Let \(h_t\) be the smooth-data Poisson solution with boundary value
\(w_t\). It is smooth up to this boundary, harmonic inside, and its
difference from \(w_t\) has zero Sobolev boundary values. For completeness,
the latter follows without an assumed trace theorem: multiply the smooth
zero-boundary difference by cutoffs vanishing in a collar of thickness
\(s\). Its magnitude is bounded by a constant times distance to the sphere;
therefore the cutoff-gradient contribution has squared integral bounded
by a constant times the collar volume, which tends to zero. The remaining
function and gradient errors tend to zero by dominated convergence. The
cutoff products are smooth and compactly supported in the ball.

Energy orthogonality for smooth data gives \eqref{eq:orthogonal} for
\(w_t,h_t\). Applied to the harmonic difference \(h_t-h_s\), it also gives
\[
 \norm{Dh_t-Dh_s}_{L^2(B_r(x))}
 \le\norm{Dw_t-Dw_s}_{L^2(B_r(x))}.
\]
The Poisson probability kernel bounds
\(\sup|h_t-h_s|\) by the boundary supremum of \(|w_t-w_s|\).
Consequently \(h_t\) converges uniformly on the closed ball and its
gradients converge strongly in \(L^2\). Passing distributional pairings
to the limit identifies these gradients with those of the limit \(h\).
The differences \(w_t-h_t\) converge in the Sobolev graph norm, so their
limit belongs to the same closed zero-boundary space. Uniform convergence
preserves the boundary value. Passing the mean-value identity to the
uniform limit gives the identity for \(h\); the local radial-kernel formula
then proves its interior smoothness and harmonicity. Strong gradient
convergence proves \eqref{eq:orthogonal}.

The signed Green identity for the smooth pair is
\[
 \int_{B_r(x)}|Dw_t-Dh_t|^2
 =-\int_{B_r(x)}F_t\cdot(w_t-h_t).
\]
The oscillation of \(w_t\) around \(w(x)\) is at most
\(H(r+t)^\alpha\); the Poisson maximum estimate therefore gives
\(\sup|w_t-h_t|\le2H(r+t)^\alpha\). It follows that the error energy is
at most
\(2KH(r+t)^\alpha\int_{B_{r+t}(x)}|Dw|^2\).
The last integral converges to the integral over \(B_r(x)\): decreasing
closed-radius neighborhoods converge to the closed ball, the sphere has
Lebesgue measure zero, and the energy is integrable in the larger ball.
Strong gradient convergence yields \eqref{eq:error}. Thus no direct use of
an unproved bounded Sobolev test extension of the weak equation is needed.
\end{proof}

\begin{lemma}[Energy and gradient-excess comparison]
\label{lem:comparisons}
Write \(E_x(r)=\int_{B_r(x)}|Dw|^2\) and
\(\Phi_x(r)=\int_{B_r(x)}|Dw-(Dw)_{x,r}|^2\). For
\(0<\theta\le1/2\),
\begin{align}
 E_x(\theta r)&\le(c_n\theta^n+cKHr^\alpha)E_x(r),
 \label{eq:energycomparison}\\
 \Phi_x(\theta r)&\le c_n\theta^{n+2}\Phi_x(r)
                  +cKHr^\alpha E_x(r).
 \label{eq:excesscomparison}
\end{align}
The constants are independent of the center, radius, map, and energy.
\end{lemma}
\begin{proof}
Decompose \(Dw=Dh+(Dw-Dh)\), use the squared triangle inequality,
\eqref{eq:henergy}, and \eqref{eq:orthogonal} for the first inequality.
For the second, the ball mean minimizes the integral squared distance
from a constant matrix. Compare \(Dw\) to \((Dh)_{x,\theta r}\), apply
\eqref{eq:hexcess}, and observe
\[
 \int_{B_r(x)}|Dh-(Dh)_{x,r}|^2
 \le2\Phi_x(r)+2\int_{B_r(x)}|Dw-Dh|^2.
\]
The last inequality follows by comparison with \((Dw)_{x,r}\).
Insert \eqref{eq:error}; since \(\theta^{n+2}\le1\), all error terms
have the stated bound. This derivation does not assume classical
differentiability of \(w\).
\end{proof}

\section{Improved energy growth and uniform parameter selection}
\label{sec:growth}

Set
\begin{equation}\label{eq:parameters}
 r_0=\frac1{16},\qquad \beta=1-\frac\alpha4,\qquad
 \lambda=n-2+2\beta=n-\frac\alpha2,\qquad
 \gamma=\frac\alpha8.
\end{equation}
For every \(x\in B_{5/8}\) and \(r\le r_0\), the replacement ball is
compactly contained in \(B_{3/4}\). Enlarge the comparison constants once
to a number \(c\ge1\). Choose \(\theta_1\in(0,1/4]\) such that
\(c\theta_1^{n-\lambda}\le1/4\), possible because
\(n-\lambda=\alpha/2>0\). Choose \(H_0>0\) so that
\begin{equation}\label{eq:smallholder}
 c(1+K)H_0r_0^\alpha\le\tfrac14\theta_1^\lambda.
\end{equation}
The harmless \(1+K\) also covers \(K=0\).

\begin{lemma}[Uniform improved growth]
\label{lem:growth}
If \(H\le H_0\), then for all the centers and radii above,
\begin{equation}\label{eq:growth}
 E_x(r)\le C_E E\,r^\lambda.
\end{equation}
\end{lemma}
\begin{proof}
The comparison inequality gives
\(E_x(\theta_1r)\le\tfrac12\theta_1^\lambda E_x(r)
\le\theta_1^\lambda E_x(r)\).
For \(r_j=\theta_1^jr_0\), induction gives
\(E_x(r_j)\le(r_j/r_0)^\lambda E_x(r_0)\).
If \(r_{j+1}<r\le r_j\), monotonicity of the ball integral gives
\(E_x(r)\le E_x(r_j)\le\theta_1^{-\lambda}(r/r_0)^\lambda E\).
This proves the bound with
\(C_E=\theta_1^{-\lambda}r_0^{-\lambda}\), including every radius.
No division by \(E\) occurs, so the argument includes \(E=0\).
\end{proof}

Combining \eqref{eq:error} and \eqref{eq:growth} gives
\begin{equation}\label{eq:bettererror}
 \int_{B_r(x)}|Dw-Dh|^2
 \le C H E\,r^{\lambda+\alpha}
 \le C E\,r^{n+\alpha/2}.
\end{equation}
The improvement beyond the exponent \(n\) is the reason for selecting
\(\beta\) close to one before treating the gradient.

Select the final energy threshold now:
\begin{equation}\label{eq:threshold}
 \varepsilon_*=
 \min\{1,\ 4^{2-n}\varepsilon_H,\ (H_0/L)^2\}>0.
\end{equation}
The constant \(L>0\) may be enlarged if necessary. Under this threshold,
Lemma~\ref{lem:localholder} supplies \(H\le L\sqrt E\le H_0\), and
\(\sqrt E\le1\). No later derivative order changes this threshold.

\section{Gradient oscillation, Campanato representatives, and classical derivatives}
\label{sec:c1}

\begin{lemma}[Gradient excess decay]
\label{lem:excess}
For \(x\in B_{5/8}\), \(0<r\le r_0\),
\begin{equation}\label{eq:excess}
 \Phi_x(r)\le C_\Phi E\,r^{n+2\gamma}.
\end{equation}
\end{lemma}
\begin{proof}
By \eqref{eq:excesscomparison} and \eqref{eq:bettererror},
\(\Phi_x(\theta r)\le c\theta^{n+2}\Phi_x(r)
 +CEr^{n+\alpha/2}\).
Put \(\nu=n+2\gamma=n+\alpha/4\),
\(\mu=n+\alpha/2>\nu\), and choose \(\theta_2\in(0,1/4]\) with
\(c\theta_2^{n+2}\le\tfrac12\theta_2^\nu\).
For \(s_j=\theta_2^jr_0\) and
\(Z_j=s_j^{-\nu}\Phi_x(s_j)\),
\[
 Z_{j+1}\le\tfrac12 Z_j+CE\theta_2^{-\nu}s_j^{\mu-\nu}.
\]
Iterating and bounding \(s_j^{\mu-\nu}\le r_0^{\mu-\nu}\) bounds
\(Z_j\) by \(Z_0+2CE\theta_2^{-\nu}r_0^{\mu-\nu}\).
Also \(\Phi_x(r_0)\le E\). To pass to an arbitrary radius
\(s_{j+1}<r\le s_j\), compare to the constant mean on \(B_{s_j}(x)\);
minimality of the mean gives \(\Phi_x(r)\le\Phi_x(s_j)\).
This costs only the factor \(\theta_2^{-\nu}\).
\end{proof}

\begin{lemma}[Campanato reconstruction]
\label{lem:campanato}
Let \(g\) be a locally square-integrable Euclidean-vector-valued field,
with \(\int_{B_{3/4}}|g|^2\le E\). Suppose
\(\int_{B_r(x)}|g-g_{x,r}|^2\le M E r^{n+2\gamma}\) for
\(x\in B_{5/8}\), \(r\le r_0\), where \(0<\gamma<1\).
There is a continuous representative \(G\) on \(B_{5/8}\) with
\begin{equation}\label{eq:campanatobound}
 \sup_{B_{5/8}}|G|+[G]_{\gamma;B_{5/8}}
 \le C(n,\gamma,r_0,M)\sqrt E.
\end{equation}
\end{lemma}
\begin{proof}
Cauchy--Schwarz and the volume scaling of balls imply
\[
 |g_{x,r/2}-g_{x,r}|
 \le C_n r^{-n/2}
 \left(\int_{B_r(x)}|g-g_{x,r}|^2\right)^{1/2}
 \le C_n\sqrt{ME}\,r^\gamma.
\]
The geometric sum \(\sum_j2^{-j\gamma}\) converges, uniformly in \(x\).
The dyadic means therefore converge uniformly to \(G(x)\), with
\(|G(x)-g_{x,r}|\le C\sqrt{ME}\,r^\gamma\).
For radii between consecutive dyadic radii the same bound follows by
comparison of means, with a volume ratio at most \(2^n\). Thus the limit
is an all-radius limit. Lebesgue differentiation identifies it with \(g\)
almost everywhere.
If \(d=|x-y|\le r_0/4\), compare the means on \(B_{2d}(x)\) and
\(B_{2d}(y)\) to that on \(B_{4d}(x)\). Both smaller balls lie in the
last ball; Cauchy--Schwarz and the oscillation bound give differences
at most \(C\sqrt{ME}\,d^\gamma\). Add the two limit-to-mean errors.
At radius \(r_0\), the mean is bounded by
\(C_nr_0^{-n/2}\sqrt E\); adding the limit error bounds \(G\) uniformly.
For \(d>r_0/4\), the uniform bound gives the desired H\"older estimate
with the additional constant \(2(r_0/4)^{-\gamma}\). These estimates prove
continuity, so the almost-everywhere representative just obtained has the
claimed pointwise meaning.
\end{proof}

Apply this lemma to the gradient matrix, viewed as a vector in its
Hilbert--Schmidt Euclidean space. Denote its continuous representative
by \(G\). Let \(T(x):\R^n\to\R^q\) have columns \(G_i(x)\).
The operator norm of \(T\) is at most the Hilbert--Schmidt norm of \(G\).
On a compact inner ball, mollification satisfies
\(w_t\to w\) and \(Dw_t\to T\) uniformly: the weak derivatives commute
with convolution, \(Du=G\) almost everywhere, and \(G\) is continuous.
For a segment in that ball, the fundamental theorem of calculus gives,
after passage to the uniform limit,
\begin{equation}\label{eq:segment}
 w(y)-w(x)=\int_0^1 T(x+s(y-x))(y-x)\dd s.
\end{equation}
Subtracting \(T(x)(y-x)\), the remainder is bounded by
\([T]_\gamma |y-x|^{1+\gamma}/(1+\gamma)\). This proves classical
differentiability, \(Dw=T\), and
\begin{equation}\label{eq:c1}
 \sup_{B_{5/8}}\norm{Dw}+[Dw]_{\gamma;B_{5/8}}
 \le C_1\sqrt E.
\end{equation}
The same continuous representative is used throughout. The weak equation
now has its classical gradient, but a pointwise second derivative has not
yet been asserted.

\begin{theorem}[Unit-ball quadratic Poisson first-derivative estimate]
\label{thm:c1gate}
Fix \(n\ge2,q\), \(\alpha\in(0,1)\), and \(K\ge0\). There are
\(H_* >0\) and \(C<\infty\), depending only on \(n,q,\alpha,K\),
such that the following holds. Let \(u\in W^{1,2}(B_1,\R^q)\) have
\([u]_{\alpha;B_1}\le H\le H_*\), \(H\ge0\), and let
\(\Delta u=F\) weakly with \(|F|\le K|Du|^2\) almost everywhere.
Then \(u\in C^{1,\alpha/8}(B_{1/2})\), and
\[
\sup_{B_{1/2}}\norm{Du}\le C\left(\int_{B_1}|Du|^2\right)^{1/2},
\qquad [Du]_{\alpha/8;B_{1/2}}
\le C\left(\int_{B_1}|Du|^2\right)^{1/2}.
\]
The original continuous function is the resulting differentiable function.
Their sum is bounded by twice the displayed constant times the energy square root.
\end{theorem}
\begin{proof}
The arguments of Sections~\ref{sec:comparison}--\ref{sec:c1} use only
the quadratic bound, H\"older continuity, and finite Sobolev energy.
Apply them with centers in \(B_{1/2}\), starting radius \(1/16\),
and sampling/replacement balls compactly contained in \(B_1\).
Select \(H_*\) by \eqref{eq:smallholder}. All limit arguments are local
on a slightly larger ball inside \(B_1\); the initial energy is bounded
by the whole-ball energy. Campanato reconstruction and
\eqref{eq:segment} give exactly the asserted representative and estimates.
Increasing \(C\) once if needed gives both displayed estimates. No target
constraint or minimality is used in this theorem.
\end{proof}

\begin{corollary}[Local qualitative upgrade]
\label{cor:qualitative}
A H\"older-continuous weak solution of
\(\Delta v=\sum_i A(v)[D_iv,D_iv]\), with smooth \(A\), is locally smooth.
More precisely, \(A\) is a smooth bilinear-map-valued function on
\(\R^q\), \(v\in W^{1,2}_{\mathrm{loc}}(\Omega,\R^q)\), and every
point has a compactly contained ball on which \(v\) has a finite H\"older
bound with an exponent in \((0,1)\). The exponent and bound may vary
between these local balls. The weak equation includes its locally
integrable forcing and the signed test identity.
No minimality is needed in this upgrade assertion.
\end{corollary}
\begin{proof}
Fix a point and a Holder ball about it, then restrict to a concentric
closed half-ball. Holder continuity implies continuity there, so its
image is compact and the coefficient and all its derivatives have finite
bounds on a compact ambient neighborhood of that image. Let \(K\ge0\)
bound the resulting quadratic forcing coefficient. Choose a smaller
radius \(\rho>0\), still compactly contained, such that
\(H\rho^\alpha\le H_*\), where \(H_*\) is the threshold of
Theorem~\ref{thm:c1gate}. This is possible even for \(H=0\), for instance
by taking \(\rho\le(H_*/(1+H))^{1/\alpha}\) in addition to the containment
restriction. For \(v_\rho(z)=v(b+\rho z)\), the actual weak gradient is
\(\rho Dv(b+\rho z)\), its Holder bound is \(H\rho^\alpha\), and its
weak forcing is \(F_\rho(z)=\rho^2F(b+\rho z)\). To verify the last
identity, test the original equation with
\(\psi((x-b)/\rho)\). Its derivative has the factor \(\rho^{-1}\),
and changing measure contributes \(\rho^n\); multiplication by
\(\rho^{2-n}\) gives exactly
\(\int Dv_\rho:D\psi=-\int F_\rho\cdot\psi\).
The quadratic coefficient bound \(K\) is unchanged. Compact-local
Sobolev membership gives finite whole-ball Sobolev data after this
rescaling. Theorem~\ref{thm:c1gate} gives local \(C^{1,\gamma}\)
regularity. The next two sections give local higher regularity by repeated
application of the proved Poisson estimate. For this qualitative application,
take the parameter \(\delta=1\) in the composition estimate and let its
constant include the finite local derivative bounds already obtained.
This qualitative assertion
does not claim constants uniform over maps with unrestricted H\"older norm.
\end{proof}

\section{A proved constant-coefficient Poisson Schauder estimate}
\label{sec:schauder}

\begin{lemma}[Interior Poisson estimate]
\label{lem:schauder}
Let \(0<s<t\), \(0<\eta<1\). Suppose \(v\) is continuous and locally
Sobolev on \(B_t\), \(f\in C^{0,\eta}(B_t,\R^q)\) with finite norm,
and \(\Delta v=f\) weakly. Then \(v\in C^{2,\eta}(B_s)\), the same
equation holds classically on \(B_s\), and
\begin{equation}\label{eq:schauder}
 \norm{v-c}_{C^{2,\eta}(B_s)}
 \le C(n,q,\eta,s,t)
 \bigl(\sup_{B_t}|v-c|+\norm{f}_{C^{0,\eta}(B_t)}\bigr)
\end{equation}
for every constant vector \(c\), when the quantities on the right are finite.
The translated and scaled versions follow by the usual affine change of
variables. Only the scalar Laplacian, applied componentwise, is involved.
\end{lemma}
\begin{proof}
Choose \(s<t_1<t_2<t\) and a smooth cutoff \(\chi\), equal to one on
\(B_{t_1}\), supported in \(B_{t_2}\). Extend \(g=\chi f\) by zero.
Then \(g\) is globally H\"older and compactly supported, and
\begin{equation}\label{eq:cutoffholder}
 \norm{g}_\infty+[g]_\eta
 \le C(\eta,t_1,t_2,t)\norm{f}_{C^{0,\eta}(B_t)}.
\end{equation}
Inside \(B_t\) this is the product estimate, using the Lipschitz bound of
\(\chi\). For one point outside \(B_t\) and the other in \(\supp\chi\),
their distance is at least \(t-t_2\), which gives the same bound.

Let \(\sigma_{n-1}\) be the area of the unit sphere and, for \(z\ne0\),
\begin{equation}\label{eq:fundamental}
 \Gamma_n(z)=
 \begin{cases}
 ((2-n)\sigma_{n-1})^{-1}|z|^{2-n},&n\ge3,\\
 (2\pi)^{-1}\log|z|,&n=2.
 \end{cases}
\end{equation}
This is locally integrable, as is its first derivative. Integrating by
parts outside a ball of radius \(\epsilon\), its radial flux through
the sphere is one. The terms involving \(\Gamma_n\) itself vanish as
\(\epsilon\to0\), since their absolute bounds are \(O(\epsilon)\) for
\(n\ge3\) and \(O(\epsilon|\log\epsilon|)\) for \(n=2\). Consequently
\(\Delta\Gamma_n=\delta_0\) in distributions, with the displayed positive
Laplacian convention.

For \(z\ne0\), direct differentiation gives
\begin{equation}\label{eq:kernel}
 K_{ij}(z)=D_{ij}\Gamma_n(z)
 =\sigma_{n-1}^{-1}
 \bigl(\delta_{ij}|z|^{-n}-nz_iz_j|z|^{-n-2}\bigr).
\end{equation}
Thus \(|K_{ij}(z)|\le C_n|z|^{-n}\),
\(|DK_{ij}(z)|\le C_n|z|^{-n-1}\), and
\begin{equation}\label{eq:cancellation}
 \int_{a<|z|<b}K_{ij}(z)\dz=0\quad(0<a<b).
\end{equation}
Indeed reflection kills off-diagonal sphere integrals, coordinate
symmetry makes every diagonal integral equal, and their sum is the sphere
area. Hence \(\int\theta_i\theta_j\dd S=\sigma_{n-1}\delta_{ij}/n\),
which proves cancellation. A second integration by parts across the
deleted sphere gives the distributional identity
\begin{equation}\label{eq:deltaHessian}
 D_{ij}\Gamma_n=\PV K_{ij}+\frac{\delta_{ij}}n\delta_0.
\end{equation}
The boundary term is
\(\sigma_{n-1}^{-1}\int\theta_i\theta_j\varphi(\epsilon\theta)\dd S\),
whose limit is \(\delta_{ij}\varphi(0)/n\). This term must not be omitted.

Define \(P=\Gamma_n*g\). On any fixed bounded set its value and first
derivative are bounded by \(C\norm{g}_\infty\), since their kernels are
locally integrable and \(g\) has fixed bounded support. Its second
derivative is
\begin{equation}\label{eq:potentialHessian}
 D_{ij}P(x)=\frac{\delta_{ij}}n g(x)
 +\int_{|z|<1}K_{ij}(z)\bigl(g(x-z)-g(x)\bigr)\dz
 +\int_{|z|\ge1}K_{ij}(z)g(x-z)\dz.
\end{equation}
The near integral is absolutely convergent, bounded by
\(C[g]_\eta\int_0^1r^{\eta-1}\dd r=C[g]_\eta/\eta\).
The far integral, on the bounded set in question, is bounded by a constant
times \(\norm{g}_\infty\). These bounds are uniform in \(x\).

Here are the details of the H\"older estimate, where cancellation is
essential. Let \(d=|x-y|\), \(d>0\), and let \(D=B_{2d}(x)\).
Write each principal value as the limit over expanding balls centered
at its evaluation point, and subtract its constant value of \(g\), using
\eqref{eq:cancellation}. On \(D\), the two integrals with differences
\(g(z)-g(x)\) and \(g(z)-g(y)\) are bounded by \(C[g]_\eta d^\eta\):
the relevant distances are at most \(2d\) and \(3d\). The extra constant
term from changing the subtracted value is also bounded by this quantity,
because
\[
 B_d(y)\subset D\subset B_{3d}(y),\qquad
 \left|\PV\int_D K_{ij}(y-z)\dz\right|
 \le C\int_d^{3d}\frac{\dd r}{r}=C\log3.
\]
On \(D^c\), the mean-value bound for the kernel gives
\[
 |K_{ij}(x-z)-K_{ij}(y-z)|\le Cd|x-z|^{-n-1}.
\]
Multiplying by \(|g(z)-g(x)|\le[g]_\eta|x-z|^\eta\) and integrating
gives
\(C[g]_\eta d\int_{2d}^\infty r^{\eta-2}\dd r
 \le C[g]_\eta d^\eta/(1-\eta)\).
The remaining constant term on \(D^c\) is the negative of the integral
on \(D\), since the full centered annular integrals vanish; it is bounded
by the same \(\log3\) estimate. When expanding the outer balls, the
translation difference of their boundaries contributes \(O(d/R)\)
times a bounded value of \(g\), and vanishes as \(R\to\infty\).
The explicit \(g(x)\) term in \eqref{eq:potentialHessian} is H\"older as
well. This proves the required full exponent \(\eta\) estimate, without
a logarithmic loss.

To justify classical differentiability for a merely H\"older \(g\),
mollify \(g\) to \(g_\epsilon\). The potentials of these smooth data are
smooth: differentiate the compactly supported data factor in
\(\int\Gamma_n(z)g_\epsilon(x-z)\dz\); on every compact set of \(x\)
the locally integrable kernel gives a common dominating function for each
finite number of derivatives. The supports remain in a fixed compact
set, their H\"older constants are uniformly bounded, and they converge
uniformly to \(g\). Potentials and first derivatives converge uniformly
on compact sets by local kernel integrability. For the Hessian formula,
split at a fixed small radius \(a\). The near contribution to the
difference is bounded by \(C[g]_\eta a^\eta\), uniformly in
\(\epsilon\); on the complement uniform convergence of \(g_\epsilon\)
gives convergence of the integrals. First let \(\epsilon\to0\), then
\(a\to0\). Thus the Hessians converge uniformly to
\eqref{eq:potentialHessian}. The fundamental theorem of calculus for
the smooth potentials and their first derivatives identifies both limit
derivatives. In particular \(P\in C^{2,\eta}\) locally. Taking the trace
in \eqref{eq:potentialHessian} gives \(\Delta P=g\), because
\(\sum_iK_{ii}=0\) and \(\sum_i\delta_{ii}/n=1\).

Now \(h=v-P\) is continuous and weakly harmonic on \(B_{t_1}\).
Its mollifications are harmonic on slightly smaller balls and converge
uniformly on compact sets. Passing their mean-value identities to the
limit and using the radial-kernel formula proves that \(h\) is smooth
and harmonic. Its interior derivatives of orders up to three are bounded
by a constant times \(\sup_{B_{t_1}}|h-c|\), as in
\eqref{eq:hderivatives}. The third derivative bounds the H\"older
seminorm of the second derivative on \(B_s\), since the ball is convex
and \(|x-y|\le2s\). Combine these estimates for \(h\) with those for
\(P\), and use \eqref{eq:cutoffholder}. This proves \eqref{eq:schauder}.
\end{proof}

\begin{corollary}[Unit-ball pointwise Schauder interface]
\label{cor:unitschauder}
For fixed \(n\ge2,q\) and \(\eta\in(0,1)\), there is \(C\ge0\),
chosen before the solution and bounds, such that the following holds.
Let \(v\in W^{1,2}_{\mathrm{loc}}(B_1,\R^q)\cap C(B_1,\R^q)\)
weakly solve \(\Delta v=f\). Suppose
\(\sup_{B_1}|v-c|\le M\), \(\sup_{B_1}|f|\le L\), and
\([f]_{\eta;B_1}\le H\), with \(M,L,H\ge0\). Then
\(v\in C^{2,\eta}(B_{1/2})\), the equation holds pointwise, and
\[
 \sup_{B_{1/2}}\norm{D^j(v-c)}\le C(M+L+H)\quad(0\le j\le2),
 \qquad [D^2v]_{\eta;B_{1/2}}\le C(M+L+H).
\]
Here \(D^0(v-c)=v-c\).
\end{corollary}
\begin{proof}
Use Lemma~\ref{lem:schauder} with \(s=1/2,t=1\). Conversely this unit
interface supplies any fixed inner-radius version needed in the induction:
put \(d=t-s\) and apply it on radius-\(d/2\) balls centered in \(B_s\).
The rescaled forcing is \((d/2)^2f\), with H\"older bound
\((d/2)^{2+\eta}H\). This gives uniform derivative bounds at every
center. For pairs at distance below \(d/8\), use the same local estimate;
for larger distances use twice the derivative supremum divided by
\((d/8)^\eta\). Thus the global H\"older estimate on \(B_s\) also has
a constant depending only on \(n,q,\eta,s,t\).
\end{proof}

\section{Higher-order estimates, one threshold, and affine transport}
\label{sec:higher}

The extension \(Q\) in \eqref{eq:coefficient} is globally smooth.
At this point replace the forcing by its actual continuous representative
\begin{equation}\label{eq:forcingrepresentative}
 F_{\mathrm{cl}}(x)=\sum_i A(w(x))[D_iw(x),D_iw(x)],
\end{equation}
where the derivatives are now classical. It agrees almost everywhere with
the original \(L^1\) forcing, because the classical and displayed weak
gradients agree almost everywhere. Thus all weak test identities are
unchanged. In this section \(F\) denotes \(F_{\mathrm{cl}}\). Its pointwise
Holder and higher derivative estimates follow from the formula; no pointwise
regularity is asserted for an arbitrary representative of the original forcing.
Choose a closed Euclidean ball containing \(N\) in its interior. Every
derivative \(D^jA\) has a finite supremum on this ball, by continuity and
compactness. The ball is convex, so one additional derivative bounds the
Lipschitz variation of \(D^jA\) between target values. Denote the resulting
coefficient bounds by \(a_j\); they depend only on the fixed target and
the chosen extension, which was selected before the map.

\begin{lemma}[Quadratic composition estimates]
\label{lem:nonlinear}
Suppose \(w\in C^{m,\gamma}(B_s,N)\), \(m\ge1\), and for
\(1\le j\le m\) its derivative suprema and highest H\"older seminorm
are at most \(M\delta\), with \(0\le\delta\le1\). Then
\begin{equation}\label{eq:nonlinear}
 \left\|\sum_i A(w)[D_iw,D_iw]\right\|_{C^{m-1,\gamma}(B_s)}
 \le P_m(M,a_0,\ldots,a_m,n,q,s)\delta^2,
\end{equation}
for a finite function of the displayed data. Lower derivative H\"older
seminorms can be included in \(M\), by the next derivative bound on a
convex ball. The expression is zero if \(\delta=0\).
\end{lemma}
\begin{proof}
Repeated product and chain rules express each coordinate derivative of
order \(\ell\le m-1\) as a finite sum of contractions of
\[
 (D^rA)(w),\quad D^{b_1}w,\ldots,D^{b_r}w,
 \quad D^{c_1+1}w,\ D^{c_2+1}w,
 \qquad \sum b_i+c_1+c_2=\ell,
\]
where \(b_i\ge1\) and \(c_1,c_2\ge0\). Thus every term contains at
least two positive-order derivatives of \(w\), all of order at most
\(m\). Its supremum is bounded by a constant times
\(M^{r+2}\delta^{r+2}\le M^{r+2}\delta^2\).
For the H\"older bound, subtract the values of a product at two points
by replacing one factor at a time. For a derivative of \(w\) the
corresponding H\"older factor is bounded by \(M\delta|x-y|^\gamma\).
For \((D^rA)(w)\), use its Lipschitz bound and
\([w]_\gamma\le C(s)\sup|Dw|\le C(s)M\delta\).
Each term still contains at least the two derivative factors, and therefore
has a factor \(\delta^2\). There are finitely many terms depending only
on \(m,n,q\); their sum defines \(P_m\). This also explains precisely
which higher coefficient bounds are required.
\end{proof}

Put \(\delta=\sqrt E\) and
\begin{equation}\label{eq:nested}
 s_m=\frac12+\frac18\,2^{-(m-1)},\qquad m\ge1.
\end{equation}
Thus \(s_1=5/8\), \(s_m>s_{m+1}>1/2\), and the limit is \(1/2\).
By \eqref{eq:c1} and the H\"older oscillation of \(w\),
\(\|w-w(0)\|_{C^{1,\gamma}(B_{s_1})}\le M_1\delta\).
Suppose this assertion has been proved with order \(m\) and ball
\(B_{s_m}\), enlarging \(M_m\) to include all lower-order H\"older norms.
Lemma~\ref{lem:nonlinear} bounds \(F\) in \(C^{m-1,\gamma}(B_{s_m})\)
by \(P_m\delta^2\).

For \(m=1\), apply Lemma~\ref{lem:schauder} to \(w-w(0)\).
For \(m\ge2\), take every coordinate derivative \(\partial^I w\)
of order \(|I|=m-1\). Distributional derivatives commute, so
\(\Delta\partial^Iw=\partial^IF\) weakly on \(B_{s_m}\).
This identity follows by testing the original equation with
\((-1)^{|I|}\partial^I\psi\); its support is still compactly contained,
and integration by parts identifies the classical derivatives already
available at this order. These maps are at least \(C^{1,\gamma}\).
To verify their native Sobolev data explicitly, any \(C^1\) map on an
open Euclidean set and its first derivatives are bounded on each compact
subset. On a slightly larger compact neighborhood multiply by a smooth
cutoff equal to one near that subset, then mollify. The approximations
converge with their first derivatives uniformly there: differentiation
under the integral and uniform continuity of the original \(C^1\)
derivative give this convergence. Pass the smooth distributional
gradient test identities to the uniform limit. Finite measure of compact
sets also gives the required local \(L^2\) membership. Thus
\begin{equation}\label{eq:c1sobolev}
 C^1_{\mathrm{loc}}\ \Longrightarrow\
 W^{1,2}_{\mathrm{loc}}\text{ with the actual classical derivative},
\end{equation}
which supplies the locally Sobolev inputs to the Poisson estimate.
Applying that estimate between \(s_m\) and \(s_{m+1}\), and combining
the finite coordinate derivatives, gives
\[
 \|w-w(0)\|_{C^{m+1,\gamma}(B_{s_{m+1}})}
 \le S_m(M_m\delta+P_m\delta^2)
 \le S_m(M_m+P_m)\delta.
\]
Here \(S_m\) includes the coordinate/operator-norm equivalence constants
and the two radii. Define \(M_{m+1}\) by this finite recursion, increasing
it to include lower-order seminorms. All choices depend only on \(n,N,m\),
not on the map or \(E\). The fixed condition \(\delta\le1\) was already
ensured by \eqref{eq:threshold}.

For every finite \(m\), restriction to \(B_{1/2}\) gives the same
\(C^{m,\gamma}\) estimate. These are derivatives of one function,
not a sequence of independently selected representatives. Therefore
\(w\in C^\infty(B_{1/2},N)\), and \eqref{eq:main} holds at unit scale.
If \(E=0\), \eqref{eq:c1} gives \(Dw=0\); the segment identity makes
\(w\) constant, and all estimates have zero right side.

Finally let \(v(z)=u(a+Rz)\), with weak gradient
\(Dv(z)=RDu(a+Rz)\). Affine change of variables gives
\(\E(v;B_1)=\Th_u(a,R)\). The established Sobolev and minimizer
transport shows that \(v\) satisfies the same hypotheses, so the
unit-ball theorem applies with the same constants. Transport its
representative back by \(z=(x-a)/R\). Almost-everywhere equality follows
from the affine measure change, and it agrees with the original H\"older
representative by continuity and almost-everywhere uniqueness. Repeated
affine differentiation gives
\[
 D^jw(x)=R^{-j}D^j\widetilde v((x-a)/R),\qquad
 [D^jw]_{\gamma;B_{R/2}(a)}
 =R^{-j-\gamma}[D^j\widetilde v]_{\gamma;B_{1/2}}.
\]
This proves Theorem~\ref{thm:main}, including \eqref{eq:zero}, the
unchanged representative, and the one-threshold quantifier order.

\paragraph{Scope of the complete small-energy theorem.}
The H\"older input plus this upgrade yields smoothness and quantitative
ambient derivatives on the half-ball. Singular-set dimension estimates,
boundary estimates, existence of minimizers, and Riemannian source-domain
extensions are separate mathematical results and are not conclusions of
Theorem~\ref{thm:main}.

\begin{thebibliography}{9}
\bibitem{GM} Mariano Giaquinta and Luca Martinazzi,
\emph{An Introduction to the Regularity Theory for Elliptic Systems,
Harmonic Maps and Minimal Graphs}, second edition, 2012.
The harmonic comparison route is Theorem 9.7, pp. 215--216;
the continuous weak-map consequence is Proposition 10.40, p. 265.
Chapter 5 gives the Campanato and Schauder background.
\bibitem{SU} Richard Schoen and Karen Uhlenbeck,
\emph{A regularity theory for harmonic maps}, Journal of Differential
Geometry \textbf{17} (1982), 307--335.
\href{https://doi.org/10.4310/jdg/1214436923}{doi:10.4310/jdg/1214436923}.
\bibitem{Holder} \emph{Schoen--Uhlenbeck small-energy H\"older regularity},
Lean development and accompanying mathematical proof,
\href{https://github.com/BrookWW/SchoenUhlenbeckSmallEnergyRegularityFormalization/tree/2a46c1e82d0580da7e1c9f62c58d0349c911cc8e}{pinned source revision 2a46c1e}.
The reused inputs are \texttt{smallEnergyRegularity},
\texttt{outerVariation\_weakEquation}, \texttt{tubularProjection\_complete},
and the focused Sobolev, Poisson, Green, and affine transport modules.
\end{thebibliography}
\end{document}
